% ============================================================
%  Práctica 2 — MAAAI · memoria.tex (fichero principal)
%  Cada sección está en memoria/secciones/ y se incluye con \input.
%  No editar este fichero para escribir contenido: hacerlo en su sección.
% ============================================================
\documentclass[11pt,a4paper]{article}

\usepackage[utf8]{inputenc}
\usepackage[T1]{fontenc}
\usepackage[spanish,es-tabla]{babel}   % es-tabla: "Tabla" en vez de "Cuadro"
\usepackage{lmodern}
\usepackage[margin=2.5cm]{geometry}
\usepackage{amsmath,amssymb}
\usepackage{graphicx}
\usepackage{booktabs}
\usepackage{float}
\usepackage{caption}
\usepackage[hidelinks]{hyperref}

\graphicspath{{../results/}}           % las figuras se guardan en results/

\begin{document}

% ---------------- Portada ----------------
\begin{titlepage}
\centering
{\large Universidade da Coruña\par}
{\large Grado en Inteligencia Artificial\par}
\vspace{1cm}
{\Large Modelos Avanzados de Aprendizaje Automático I\par}
\vspace{2cm}
{\huge\bfseries Práctica 2\par}
\vspace{0.5cm}
{\LARGE Clasificación de gases con un array de sensores\\ sometido a deriva temporal\par}
\vfill
{\large
Oliver Konrad Martínez Malec\\
Pedro Otero Pago\\
Lucas Vilariño Lorenzen\par}
\vspace{1cm}
{\large Curso 2026/27\par}
\end{titlepage}

\tableofcontents
\newpage

% ---------------- Secciones (orden del enunciado) ----------------
% Memoria · Introducción y descripción del dataset (Oliver)
\section{Introducción y descripción del dataset}
\label{sec:dataset}

Pendiente.
        % 1. Introducción y descripción del dataset   (Oliver)
% Memoria · Análisis exploratorio inicial (Oliver)
\section{Análisis exploratorio inicial}
\label{sec:exploratorio}

Pendiente.

\subsection{Criterios de limpieza}
\label{sec:limpieza}

Pendiente.   % 2. Análisis exploratorio inicial            (Oliver)
% ============================================================
%  Memoria · Tratamiento y análisis de valores perdidos (Pedro)
% ============================================================
\section{Tratamiento y análisis de valores perdidos}
\label{sec:ausentes}

\subsection{Criterio de valor ausente}

Siguiendo el criterio indicado por el profesor, se considera valor ausente todo cero exacto
en cualquiera de las 128 variables sensoriales. El análisis se realiza sobre el conjunto de
datos original (13\,910 mediciones), antes de aplicar la limpieza por reglas fijas, con el fin
de caracterizar el fenómeno completo y no solo la parte que sobrevive a dicha limpieza.

Conviene señalar que este criterio no recoge todas las mediciones inválidas: un sensor que no
responde puede registrar una variación de resistencia negativa en lugar de nula. Ese tipo de
mediciones se identifica y se trata en la limpieza de datos (Sección~\ref{sec:limpieza}),
eliminando la medición completa, y no mediante imputación.

\subsection{Cuantificación}

Se detectan 55 celdas ausentes sobre un total de 1\,780\,480 ($13\,910 \times 128$), lo que
supone un 0,0031~\% del conjunto de datos. Los ausentes afectan únicamente a 6 de las 128
variables, y ninguna supera el 0,11~\% de sus valores (Tabla~\ref{tab:ausentesVariable}).

La distribución, sin embargo, está muy lejos de ser uniforme. De las 55 celdas, 52 corresponden
al sensor 4 y 51 se encuentran en el batch 2, mientras que los batches 1 y 4--9 no presentan
ningún valor ausente (Tabla~\ref{tab:ausentesBatch}). En total, solo 18 de las 13\,910
mediciones tienen algún valor ausente.

\begin{table}[h]
\centering
\caption{Valores ausentes por variable (ceros exactos sobre las 13\,910 mediciones originales).}
\label{tab:ausentesVariable}
\begin{tabular}{lrr}
\hline
Variable & Celdas ausentes & \% de la variable \\
\hline
S04\_DR       & 14 & 0,101 \\
S04\_EMAd0001 & 13 & 0,093 \\
S04\_EMAd01   & 13 & 0,093 \\
S04\_EMAd001  & 12 & 0,086 \\
S05\_DR       &  2 & 0,014 \\
S14\_DR       &  1 & 0,007 \\
\hline
Total (128 variables) & 55 & 0,0031 \\
\hline
\end{tabular}
\end{table}

\begin{table}[h]
\centering
\caption{Distribución de los valores ausentes por batch (los batches 1 y 4--9 no tienen ninguno).}
\label{tab:ausentesBatch}
\begin{tabular}{crrr}
\hline
Batch & Mediciones & Mediciones afectadas & Celdas ausentes \\
\hline
2  & 1\,244 & 14 & 51 \\
3  & 1\,586 &  2 &  2 \\
10 & 3\,600 &  2 &  2 \\
\hline
Total & 13\,910 & 18 & 55 \\
\hline
\end{tabular}
\end{table}

\subsection{Patrones de ausencia}

Para estudiar si los ausentes aparecen aislados o agrupados, se registra, en cada medición
afectada, el conjunto de variables que faltan a la vez. Se obtienen solo cinco patrones
distintos (Tabla~\ref{tab:patronesAusentes}).

\begin{table}[h]
\centering
\caption{Patrones de ausencia: variables que faltan simultáneamente en una misma medición.}
\label{tab:patronesAusentes}
\begin{tabular}{lr}
\hline
Variables ausentes & Mediciones \\
\hline
S04\_DR, S04\_EMAd0001, S04\_EMAd001, S04\_EMAd01 & 12 \\
S04\_DR                                          &  2 \\
S05\_DR                                          &  2 \\
S04\_EMAd0001, S04\_EMAd01                       &  1 \\
S14\_DR                                          &  1 \\
\hline
\end{tabular}
\end{table}

El patrón dominante, presente en 12 mediciones, consiste en la pérdida simultánea de la
variación de resistencia y de las tres características de decaimiento del sensor 4, mientras
que su respuesta normalizada y sus características de crecimiento se conservan. Este
comportamiento es compatible con un sensor que no completa su respuesta en esa medición, y no
con fallos de lectura aislados.

Los ausentes también se agrupan en el tiempo. Las mediciones afectadas forman bloques
consecutivos: 669--678 (batch 2, etanol a 500--600~ppmv), la ventana 825--833 (batch 2,
etileno a 250~ppmv) y 2722 y 2724 (batch 3, etanol a 10~ppmv). Estas 16 mediciones coinciden
con episodios que el análisis exploratorio ya había identificado como fallos de los sensores
(Sección~\ref{sec:limpieza}). Las dos mediciones restantes (10\,451 y 13\,481) pertenecen al
batch 10, corresponden a tolueno a 1~ppmv, la concentración más baja del conjunto, y solo
pierden la variable S05\_DR; el resto de la medición es normal.

\subsection{Mecanismo de ausencia}

\textbf{MCAR.} El mecanismo MCAR exige que la probabilidad de que falte un valor sea
independiente de cualquier otra información. Los resultados anteriores lo descartan: los
ausentes se concentran en un sensor, en un batch y en bloques temporales concretos. Con tan
pocas celdas, un contraste formal (por ejemplo, $\chi^2$ entre ausencia y batch) tendría
frecuencias esperadas demasiado bajas para ser fiable, pero la evidencia descriptiva es
suficientemente clara.

\textbf{MAR.} Las 16 mediciones del batch 2 y del batch 3 presentan ausentes asociados a fallos
del hardware en periodos concretos. Su aparición se puede anticipar a partir de información
observada: el batch, la posición temporal de la medición y el comportamiento anómalo del resto
de variables de la misma medición y del sensor gemelo. Por ejemplo, en la ventana 825--833 el
sensor 3 no responde y el sensor 4 presenta medias móviles rápidas entre 13 y 19 veces por
encima de lo habitual. Este comportamiento es el propio de un mecanismo MAR.

\textbf{MNAR.} Los dos ausentes de S05\_DR con tolueno a 1~ppmv no están asociados a ningún
fallo: el resto de la medición es normal. La explicación más plausible es que la respuesta del
sensor se sitúa en el límite de detección, de modo que el valor falta precisamente por ser muy
pequeño. En ese caso la ausencia depende del propio valor no observado, lo que es compatible
con un mecanismo MNAR. Con solo dos celdas esta hipótesis no puede confirmarse, pero es
coherente con el ruido observado en la respuesta del tolueno a esta concentración.

En conclusión, el mecanismo no es único: predomina un mecanismo MAR ligado a fallos de los
sensores, con un componente residual compatible con MNAR.

\subsection{Tratamiento e imputación}

Cada tipo de ausencia recibe un tratamiento distinto.

\textbf{Ausentes en mediciones defectuosas (MAR).} En estas mediciones el valor nulo no es un
dato que falte de forma aislada, sino el síntoma de una medición fallida en su conjunto. Imputar
esas celdas rellenaría huecos de una medición que ya es inválida en otras variables. Por ello se
elimina la medición completa mediante la regla fija de limpieza (más del 10~\% de variables
inválidas), aplicada antes de la partición holdout. Esta regla elimina 53 de las 55 celdas
ausentes.

\textbf{Ausentes residuales (MNAR).} Tras la limpieza quedan dos celdas ausentes, ambas en
S05\_DR: una en una medición de entrenamiento (10\,451) y otra en una de test (13\,481). El
marcado de los ceros como ausentes es una regla fija, por lo que se aplica antes del holdout.
La imputación, en cambio, se ajusta a los datos y debe hacerse después.

Para imputarlas se ha implementado un imputador KNN como nodo de MLJ, situado al inicio del
pipeline y antes de la normalización Min--Max. En el ajuste guarda únicamente las mediciones
completas del conjunto de entrenamiento y el rango de cada variable. Al transformar, para cada
medición con valores ausentes calcula la distancia a esas mediciones usando solo las variables
observadas, escalando cada diferencia por el rango de su variable para que DR, del orden de
$10^4$, no domine frente a variables del orden de la unidad. El valor ausente se sustituye por
la media de los $k = 5$ vecinos más cercanos.

Al estar dentro del pipeline, el imputador se ajusta solo con los datos de entrenamiento: en la
validación cruzada, con los de cada fold, y en la evaluación final, con todo el conjunto de
entrenamiento. Así se evita cualquier fuga de información desde validación o test.

Se descarta imputar con la media, la mediana o la moda de la variable, porque mezclarían gases
y concentraciones con respuestas de órdenes de magnitud distintos. También se descarta la
interpolación temporal entre mediciones consecutivas: la medición afectada de test está
rodeada de otras mediciones de test, por lo que se estaría usando información del propio
conjunto de test. Se elige un valor de $k$ pequeño para que los vecinos correspondan al mismo
gas y a una concentración similar.

\textbf{Resultado e impacto.} El valor imputado de S05\_DR en la medición 13\,481 es 120,53,
frente a una mediana de 115,2 en las mediciones de tolueno a 1~ppmv del conjunto de
entrenamiento, y unas cien veces menor que la mediana global de $|DR|$ (11\,415). El valor es
coherente con una respuesta débil, como predice la hipótesis MNAR. Es esperable un ligero sesgo
al alza, porque los vecinos son mediciones en las que el sensor sí respondió, pero su efecto es
despreciable: tras la limpieza solo se imputan 2 celdas de 1\,775\,616. Aun así, el imputador es
necesario, ya que los clasificadores utilizados no admiten valores ausentes.

Por último, si los ausentes no se hubieran tratado antes de validar, se habrían concentrado en
los folds que contienen el batch 2, introduciendo diferencias de calidad de datos entre folds
ajenas al problema de deriva que se pretende estudiar.       % 3. Tratamiento y análisis de valores perdidos (Pedro)
% Memoria · Estrategias de preprocessing aplicadas
\section{Estrategias de preprocessing aplicadas}
\label{sec:preprocesado}

% Memoria · Partición holdout temporal (Oliver)
\subsection{Partición holdout temporal}
\label{sec:holdout}

Pendiente.

% Memoria · Validación cruzada por batches (Lucas)
\subsection{Validación cruzada por batches}
\label{sec:validacion}

Pendiente.

% Memoria · Normalización Min--Max (Lucas)
\subsection{Normalización Min--Max}
\label{sec:normalizacion}

Pendiente.

   % 4. Estrategias de preprocessing aplicadas   (varios)
% Memoria · Descripción del pipeline experimental (B2)
\section{Descripción del pipeline experimental}
\label{sec:pipeline}

Pendiente.
       % 5. Descripción del pipeline experimental
% Memoria · Resultados experimentales (B2/B3)
\section{Resultados experimentales}
\label{sec:resultados}

Pendiente.
     % 6. Resultados experimentales
% Memoria · Comparativa de modelos (B3)
\section{Comparativa de modelos}
\label{sec:comparativa}

Pendiente.
    % 7. Comparativa de modelos
% Memoria · Análisis del drift temporal (B4/B5)
\section{Análisis del drift temporal}
\label{sec:deriva}

Pendiente.
         % 8. Análisis del drift temporal
% Memoria · Discusión metodológica (B5)
\section{Discusión metodológica}
\label{sec:discusion}

Pendiente.
      % 9. Discusión metodológica
% Memoria · Conclusiones finales (B5)
\section{Conclusiones finales}
\label{sec:conclusiones}

Pendiente.
   % 10. Conclusiones finales
% Memoria · Test estadístico sobre los resultados de validación cruzada (B2/B3)
\section{Test estadístico sobre los resultados de validación cruzada}
\label{sec:estadistica}

Pendiente.
    % 11. Test estadístico sobre la validación cruzada

\appendix
% Memoria · Instrucciones de ejecución y estructura del código (Oliver)
\section{Instrucciones de ejecución y estructura del código}
\label{sec:ejecucion}

Pendiente.
      % A. Instrucciones de ejecución y estructura del código

\end{document}